% GENERATED FILE. Edit canonical lessons, then run npm run generate:books.
% canonical-source-hash: fnv1a64:8f9e228eaf5cb2bf
% canonical-lessons: BN-C419-obhab, BN-C419-opochoy, BN-C419-kipte, BN-C419-bondhok, BN-C419-sombol

\chapter{Want and Waste}
\label{ch:bn-want-and-waste}
\hlchaptermodality{419}

\begin{chapteropening}
\textbf{By the end of this chapter:} \emph{I can talk about being short of money, wasting it and making it last.}

Complete the last lesson of chapter 419: I can talk about being short of money, wasting it and making it last.
\end{chapteropening}

\section[\texorpdfstring{\bn{অভাব} (ôbhāb)}{অভাব (ôbhāb)}]{\bn{অভাব} (ôbhāb) — want; lack; need}

\label{lesson:BN-C419-obhab}



\begin{quote}
Before the new one: say the Bengali for a moneylender, then the Bengali for the rest (remaining).
\end{quote}

\subsection*{You'll want to know: \bn{অভাব}}
\textbf{\bn{অভাব}} — \emph{ôbhāb} — \textquotedblleft{}want; lack; need\textquotedblright{}.

\begin{quote}
\bn{গরমে} \bn{গ্রামে} \bn{জলের} \bn{অভাব} \bn{হয়।}
\end{quote}

\begin{quote}
\emph{gôrôme grāme jôler ôbhāb hôy.}
\end{quote}

\begin{quote}
In summer the village runs short of water.
\end{quote}

Also hardship itself: \bn{অভাবের} \bn{দিন}, hard times; \bn{টাকার} \bn{অভাব}, a shortage of money.

\begin{cousinweb}
\textbf{\bn{অভাব}} (\emph{ôbhāb}) is Sanskrit \emph{abhāva}, \textquotedblleft{}non-existence, absence\textquotedblright{} (Monier-Williams), from \emph{a-}, not, and \emph{bhāva}, being; Bengali uses it for the want of what one needs.
\end{cousinweb}

\begin{grammarlens}[title={What you've built}]
One more word for getting around.
\end{grammarlens}

\subsection*{Your turn}
\begin{itemize}
\raggedright
  \item \emph{Say it:} \emph{ôbhāb}
  \item \emph{Say it:} \emph{ôbhāb}, once more
  \item \emph{Say it:} \emph{gôrôme grāme jôler ôbhāb hôy.} — In summer the village runs short of water.
  \item \emph{Recall:} say the Bengali for a moneylender, then the Bengali for the rest (remaining), then say \emph{ôbhāb} again
\end{itemize}

\begin{tcolorbox}[breakable,colback=teal!4,colframe=teal!35!black,title={Before you move on}]
What does \bn{অভাব} mean? (\textquotedblleft{}Want; lack; need\textquotedblright{}.) Say it once more.
\end{tcolorbox}
\section[\texorpdfstring{\bn{অপচয়} (ôpôchôy)}{অপচয় (ôpôchôy)}]{\bn{অপচয়} (ôpôchôy) — waste; wastefulness}

\label{lesson:BN-C419-opochoy}



\begin{quote}
Before the new one: say the Bengali for the rest (remaining), then the Bengali for want (lack).
\end{quote}

\subsection*{You'll want to know: \bn{অপচয়}}
\textbf{\bn{অপচয়}} — \emph{ôpôchôy} — \textquotedblleft{}waste; wastefulness\textquotedblright{}.

\begin{quote}
\bn{আলো} \bn{জ্বালিয়ে} \bn{রাখা} \bn{বিদ্যুতের} \bn{অপচয়।}
\end{quote}

\begin{quote}
\emph{ālo jvāliye rākhā bidyuter ôpôchôy.}
\end{quote}

\begin{quote}
Leaving the lights on is a waste of electricity.
\end{quote}

\begin{cousinweb}
\textbf{\bn{অপচয়}} (\emph{ôpôchôy}) is Sanskrit \emph{apacaya}, \textquotedblleft{}diminution, decrease\textquotedblright{} (Monier-Williams); the sense of wasting is the Bengali use.
\end{cousinweb}

\begin{grammarlens}[title={What you've built}]
One more word for getting around.
\end{grammarlens}

\subsection*{Your turn}
\begin{itemize}
\raggedright
  \item \emph{Say it:} \emph{ôpôchôy}
  \item \emph{Say it:} \emph{ôpôchôy}, once more
  \item \emph{Say it:} \emph{ālo jvāliye rākhā bidyuter ôpôchôy.} — Leaving the lights on is a waste of electricity.
  \item \emph{Recall:} say the Bengali for the rest (remaining), then the Bengali for want (lack), then say \emph{ôpôchôy} again
\end{itemize}

\begin{tcolorbox}[breakable,colback=teal!4,colframe=teal!35!black,title={Before you move on}]
What does \bn{অপচয়} mean? (\textquotedblleft{}Waste; wastefulness\textquotedblright{}.) Say it once more.
\end{tcolorbox}
\section[\texorpdfstring{\bn{কিপটে} (kipṭe)}{কিপটে (kipṭe)}]{\bn{কিপটে} (kipṭe) — stingy; a miser}

\label{lesson:BN-C419-kipte}



\begin{quote}
Before the new one: say the Bengali for want (lack), then the Bengali for waste (of something useful).
\end{quote}

\subsection*{You'll want to know: \bn{কিপটে}}
\textbf{\bn{কিপটে}} — \emph{kipṭe} — \textquotedblleft{}stingy; a miser\textquotedblright{}.

\begin{quote}
\bn{কিপটে} \bn{দোকানদার} \bn{একটা} \bn{পয়সাও} \bn{ছাড়লেন} \bn{না।}
\end{quote}

\begin{quote}
\emph{kipṭe dokāndār ekṭā pôysāo chhāṛlen nā.}
\end{quote}

\begin{quote}
The stingy shopkeeper would not give up a single paisa.
\end{quote}

\begin{cousinweb}
\textbf{\bn{কিপটে}} (\emph{kipṭe}) is an everyday, rather rude word for a stingy person; its origin is not settled.
\end{cousinweb}

\begin{grammarlens}[title={What you've built}]
One more word for getting around.
\end{grammarlens}

\subsection*{Your turn}
\begin{itemize}
\raggedright
  \item \emph{Say it:} \emph{kipṭe}
  \item \emph{Say it:} \emph{kipṭe}, once more
  \item \emph{Say it:} \emph{kipṭe dokāndār ekṭā pôysāo chhāṛlen nā.} — The stingy shopkeeper would not give up a single paisa.
  \item \emph{Recall:} say the Bengali for want (lack), then the Bengali for waste (of something useful), then say \emph{kipṭe} again
\end{itemize}

\begin{tcolorbox}[breakable,colback=teal!4,colframe=teal!35!black,title={Before you move on}]
What does \bn{কিপটে} mean? (\textquotedblleft{}Stingy; a miser\textquotedblright{}.) Say it once more.
\end{tcolorbox}
\section[\texorpdfstring{\bn{বন্ধক} (bôndhôk)}{বন্ধক (bôndhôk)}]{\bn{বন্ধক} (bôndhôk) — a pledge; a mortgage}

\label{lesson:BN-C419-bondhok}



\begin{quote}
Before the new one: say the Bengali for waste (of something useful), then the Bengali for stingy.
\end{quote}

\subsection*{You'll want to know: \bn{বন্ধক}}
\textbf{\bn{বন্ধক}} — \emph{bôndhôk} — \textquotedblleft{}a pledge; a mortgage\textquotedblright{}.

\begin{quote}
\bn{দরকারে} \bn{বাবা} \bn{জমিটা} \bn{বন্ধক} \bn{রাখলেন।}
\end{quote}

\begin{quote}
\emph{dôrkāre bābā jômiṭā bôndhôk rākhlen.}
\end{quote}

\begin{quote}
In his need my father mortgaged the land.
\end{quote}

\bn{বন্ধক} \bn{রাখা}, to pawn or to mortgage; \bn{বন্ধক} \bn{ছাড়ানো}, to redeem a pledge.

\begin{cousinweb}
\textbf{\bn{বন্ধক}} (\emph{bôndhôk}) is Sanskrit \emph{bandhaka}, \textquotedblleft{}pledging, a pledge\textquotedblright{} (Monier-Williams), from \emph{bandh}, to bind: a thing bound as a surety.
\end{cousinweb}

\begin{grammarlens}[title={What you've built}]
One more word for getting around.
\end{grammarlens}

\subsection*{Your turn}
\begin{itemize}
\raggedright
  \item \emph{Say it:} \emph{bôndhôk}
  \item \emph{Say it:} \emph{bôndhôk}, once more
  \item \emph{Say it:} \emph{dôrkāre bābā jômiṭā bôndhôk rākhlen.} — In his need my father mortgaged the land.
  \item \emph{Recall:} say the Bengali for waste (of something useful), then the Bengali for stingy, then say \emph{bôndhôk} again
\end{itemize}

\begin{tcolorbox}[breakable,colback=teal!4,colframe=teal!35!black,title={Before you move on}]
What does \bn{বন্ধক} mean? (\textquotedblleft{}A pledge; a mortgage\textquotedblright{}.) Say it once more.
\end{tcolorbox}
\section[\texorpdfstring{\bn{সম্বল} (sômbôl)}{সম্বল (sômbôl)}]{\bn{সম্বল} (sômbôl) — one's means; what one has to live on}

\label{lesson:BN-C419-sombol}



\begin{quote}
Before the new one: say the Bengali for stingy, then the Bengali for a pledge (for a loan).
\end{quote}

\subsection*{You'll want to know: \bn{সম্বল}}
\textbf{\bn{সম্বল}} — \emph{sômbôl} — \textquotedblleft{}one's means; what one has to live on\textquotedblright{}.

\begin{quote}
\bn{ওই} \bn{টাকাটাই} \bn{ছিল} \bn{আমাদের} \bn{পথের} \bn{সম্বল।}
\end{quote}

\begin{quote}
\emph{oi ṭākāṭāi chhilô āmāder pôther sômbôl.}
\end{quote}

\begin{quote}
That money was all we had for the journey.
\end{quote}

\begin{cousinweb}
\textbf{\bn{সম্বল}} (\emph{sômbôl}) is Sanskrit \emph{śambala}, also written \emph{sambala}, \textquotedblleft{}provisions for a journey\textquotedblright{} (Monier-Williams): \bn{পথের} \bn{সম্বল} keeps the old sense, and the word has grown to mean all one has to live on.
\end{cousinweb}

\begin{grammarlens}[title={What you've built}]
One more word for getting around.
\end{grammarlens}

\subsection*{Your turn}
\begin{itemize}
\raggedright
  \item \emph{Say it:} \emph{sômbôl}
  \item \emph{Say it:} \emph{sômbôl}, once more
  \item \emph{Say it:} \emph{oi ṭākāṭāi chhilô āmāder pôther sômbôl.} — That money was all we had for the journey.
  \item \emph{Recall:} say the Bengali for stingy, then the Bengali for a pledge (for a loan), then say \emph{sômbôl} again
\end{itemize}

\begin{tcolorbox}[breakable,colback=teal!4,colframe=teal!35!black,title={Before you move on}]
What does \bn{সম্বল} mean? (\textquotedblleft{}One's means; what one has to live on\textquotedblright{}.) Say it once more.
\end{tcolorbox}
